% Compatibilidad con instalaciones MiKTeX cuyo formato base es anterior a
% los paquetes Beamer/Hyperref instalados.
\providecommand{\IfFormatAtLeastTF}[3]{#3}
\providecommand{\IfFormatAtLeastT}[2]{}
\providecommand{\IfFormatAtLeastF}[2]{#2}
\providecommand{\IfDocumentMetadataTF}[2]{#2}
\providecommand{\IfDocumentMetadataT}[1]{}
\providecommand{\IfDocumentMetadataF}[1]{#1}
\makeatletter
\@namedef{ver@epstopdf-base.sty}{9999/12/31}
\makeatother

\documentclass[aspectratio=169,11pt]{beamer}

% =============================================================================
% MACROECONOMIA II | LICENCIATURA EN ECONOMIA | CIDE
% SESION 9: INFLACION, CURVAS DE PHILLIPS Y NEO-FISHERISMO
% =============================================================================

\usepackage[utf8]{inputenc}
\usepackage[english]{babel}
\uselanguage{English}
\languagepath{English}
\usepackage{graphicx}
\usepackage{amsmath,amssymb}
\IfFileExists{booktabs.sty}{
  \usepackage{booktabs}
}{
  \providecommand{\toprule}{\hline}
  \providecommand{\midrule}{\hline}
  \providecommand{\bottomrule}{\hline}
  \providecommand{\cmidrule}[1]{\hline}
}
\providecommand{\addlinespace}[1][0.5em]{}
\usepackage{xcolor}
\usepackage{ragged2e}
\usepackage{etoolbox}
\usepackage{array,tabularx}
\IfFileExists{microtype.sty}{\usepackage{microtype}}{}

\newtheorem{proposition}[theorem]{Proposición}
\newtheorem{remark}[theorem]{Observación}
\newtheorem{assumption}[theorem]{Supuesto}

\tolerance=1000
\emergencystretch=1em
\AtBeginDocument{\justifying}
\addtobeamertemplate{itemize body begin}{\justifying}{}
\addtobeamertemplate{enumerate body begin}{\justifying}{}
\addtobeamertemplate{description body begin}{\justifying}{}
\addtobeamertemplate{block begin}{}{\justifying}
\addtobeamertemplate{block alerted begin}{}{\justifying}
\addtobeamertemplate{block example begin}{}{\justifying}
\AtBeginEnvironment{column}{\justifying}

\definecolor{CIDENavy}{RGB}{3,25,54}
\definecolor{CIDEBlue}{RGB}{5,35,75}
\definecolor{CIDEInk}{RGB}{25,34,45}
\definecolor{CIDEGray}{RGB}{101,112,126}
\definecolor{CIDEGrayLight}{RGB}{243,244,246}

\setbeamercolor{normal text}{fg=CIDEInk,bg=white}
\setbeamercolor{structure}{fg=CIDEBlue}
\setbeamercolor{alerted text}{fg=CIDEBlue}
\setbeamercolor{title}{bg=CIDEBlue,fg=white}
\setbeamercolor{subtitle}{fg=CIDEInk}
\setbeamercolor{frametitle}{bg=CIDEBlue,fg=white}
\setbeamercolor{palette primary}{bg=CIDEBlue,fg=white}
\setbeamercolor{palette secondary}{bg=CIDEBlue,fg=white}
\setbeamercolor{palette tertiary}{bg=CIDENavy,fg=white}
\setbeamercolor{palette quaternary}{bg=CIDENavy,fg=white}
\setbeamercolor{section in head/foot}{bg=CIDENavy,fg=white}
\setbeamercolor{section in head/foot shaded}{bg=CIDENavy,fg=white!50}
\setbeamercolor{mini frame}{bg=white,fg=white}
\setbeamercolor{mini frame shaded}{bg=white!35,fg=white!35}
\setbeamercolor{upper separation line head}{bg=CIDENavy}
\setbeamercolor{lower separation line head}{bg=CIDEBlue}
\setbeamercolor{block title}{bg=CIDEBlue,fg=white}
\setbeamercolor{block body}{bg=CIDEGrayLight,fg=CIDEInk}
\setbeamercolor{block title alerted}{bg=CIDENavy,fg=white}
\setbeamercolor{block body alerted}{bg=CIDEGrayLight,fg=CIDEInk}
\setbeamercolor{caption name}{fg=CIDEBlue}

\setbeamerfont{author in head/foot}{size=\tiny}
\setbeamerfont{institute in head/foot}{size=\tiny}
\setbeamerfont{title in head/foot}{size=\tiny}
\setbeamersize{text margin left=0.70cm,text margin right=0.70cm}
\setlength{\leftmargini}{1.2em}
\setlength{\leftmarginii}{1.1em}
\setlength{\abovecaptionskip}{0.25em}
\setlength{\belowcaptionskip}{0pt}

\useoutertheme[subsection=false]{miniframes}
\setbeamertemplate{navigation symbols}{}
\setbeamertemplate{mini frame}[circle]
\setbeamertemplate{mini frame in current subsection}[circle]
\setbeamertemplate{section in head/foot}{\color{white}\insertsectionhead}
\setbeamertemplate{section in head/foot shaded}{\color{white!55}\insertsectionhead}
\setbeamertemplate{caption}[numbered]
\setbeamertemplate{section in toc}[circle]
\setbeamertemplate{subsection in toc}[circle]
\setbeamertemplate{itemize item}{\raisebox{0.15ex}{\tiny$\blacktriangleright$}}
\setbeamertemplate{itemize subitem}{\raisebox{0.15ex}{\tiny$\bullet$}}
\setbeamertemplate{enumerate item}[circle]

\setbeamertemplate{frametitle}{%
  \nointerlineskip
  \begin{beamercolorbox}[
    wd=\paperwidth,ht=3.05ex,dp=1.05ex,
    leftskip=0.42cm,rightskip=0.42cm
  ]{frametitle}
    \usebeamerfont{frametitle}\insertframetitle
  \end{beamercolorbox}%
}

\setbeamertemplate{footline}{%
  \leavevmode
  \vbox{%
    \offinterlineskip
    \hbox{%
      \begin{beamercolorbox}[wd=.72\paperwidth,ht=2.05ex,dp=.55ex,leftskip=.35cm]{palette secondary}
        \usebeamerfont{author in head/foot}\insertshortauthor
      \end{beamercolorbox}%
      \begin{beamercolorbox}[wd=.28\paperwidth,ht=2.05ex,dp=.55ex,rightskip=.35cm plus1fil]{palette secondary}
        \hfill\usebeamerfont{institute in head/foot}\insertshortinstitute
      \end{beamercolorbox}%
    }%
    \hbox{%
      \begin{beamercolorbox}[wd=.86\paperwidth,ht=1.85ex,dp=.50ex,leftskip=.35cm]{palette tertiary}
        \usebeamerfont{title in head/foot}\insertshorttitle
      \end{beamercolorbox}%
      \begin{beamercolorbox}[wd=.14\paperwidth,ht=1.85ex,dp=.50ex,center]{palette tertiary}
        \usebeamerfont{title in head/foot}\insertframenumber{} / \inserttotalframenumber
      \end{beamercolorbox}%
    }%
  }%
  \vskip0pt
}

\newcommand{\RutaLogoCIDE}{Figuras/General/logocide.png}
\setbeamertemplate{background}{%
  \vbox to \paperheight{%
    \vfill
    \hbox to \paperwidth{%
      \hfill
      \includegraphics[height=0.58cm,keepaspectratio]{\RutaLogoCIDE}%
      \hspace{0.18cm}%
    }%
    \vspace{0.68cm}%
  }%
}

\setbeamertemplate{title page}{%
  \vspace*{0.35cm}
  \centering
  \begin{beamercolorbox}[wd=\textwidth,sep=0.24cm,center,rounded=false]{title}
    \usebeamerfont{title}\inserttitle\par
  \end{beamercolorbox}
  \vspace{0.25cm}
  {\usebeamerfont{subtitle}\usebeamercolor[fg]{subtitle}\insertsubtitle\par}
  \vfill
  {\usebeamerfont{author}\insertauthor\par}
  \vspace{0.35cm}
  {\usebeamerfont{institute}\insertinstitute\par}
  \vfill
  {\usebeamerfont{date}\insertdate\par}
  \vspace{0.25cm}
}

\newcommand{\FigurasSesion}{Figuras/09_Inflacion Curvas Phillips Neo Fisherismo}
\newcommand{\figura}[2][0.65\textheight]{%
  \includegraphics[width=\linewidth,height=#1,keepaspectratio]{\FigurasSesion/#2}%
}
\newcommand{\cidehighlight}[1]{\textcolor{CIDEBlue}{\bfseries #1}}
\newcommand{\fuentefigura}[1]{%
  \vspace{0.03cm}\par{\raggedright\tiny\textit{Fuente:} #1\par}%
}
\newcommand{\notafuentes}[1]{\note{[Sources] #1}}
\newcolumntype{Y}{>{\raggedright\arraybackslash}X}

\hypersetup{
  colorlinks=true,
  urlcolor=blue,
  linkcolor=.,
  citecolor=.,
  filecolor=.
}

\newif\ifsectiontoc
\sectiontocfalse
\AtBeginSection[]{%
  \ifsectiontoc
    \begin{frame}{Contenido}
      \tableofcontents[currentsection,hideothersubsections]
    \end{frame}
  \fi
}

\title[Macroeconomía II]{Sesión 9 \\ Inflación y curvas de Phillips}
\subtitle{Política monetaria, expectativas y neo-fisherismo}
\author[LECO]{Arturo López}
\institute[CIDE]{Centro de Investigación y Docencia Económicas (CIDE)}
\date[Otoño 2026]{Licenciatura en Economía, Otoño 2026}

\begin{document}

\begin{frame}
  \titlepage
\end{frame}

\begin{frame}{Objetivos de la sesión}
  \begin{itemize}
    \setlength{\itemsep}{0.60em}
    \item Incorporar la inflación al modelo Nuevo-Keynesiano mediante una \cidehighlight{curva de Phillips}.
    \item Analizar el mandato dual del banco central y los dilemas entre estabilidad de precios y actividad económica.
    \item Estudiar el límite inferior de la tasa nominal y el mecanismo de la \cidehighlight{orientación futura}.
    \item Introducir expectativas racionales y derivar los resultados neo-fisherianos del modelo.
    \item Evaluar las reglas de Taylor, sus posibles problemas de indeterminación y una alternativa neo-fisheriana.
  \end{itemize}
\end{frame}

\section{Inflación y curva de Phillips}

\begin{frame}{¿Por qué importa la inflación?}
  \begin{itemize}
    \setlength{\itemsep}{0.58em}
    \item La inflación anticipada genera costos de ajuste de precios, economías de efectivo y distorsiones tributarias.
    \item La inflación no anticipada redistribuye riqueza entre acreedores y deudores y altera los rendimientos reales.
    \item Con precios rígidos, no todas las empresas reajustan al mismo tiempo: la inflación modifica \textbf{precios relativos} y distorsiona la asignación de recursos.
    \item Por ello, en el modelo Nuevo-Keynesiano la estabilidad de precios tiene valor social, incluso cuando la inflación es perfectamente anticipada.
  \end{itemize}
\end{frame}

\begin{frame}{Del modelo de precios rígidos a la inflación}
  \begin{itemize}
    \setlength{\itemsep}{0.62em}
    \item En el capítulo anterior, algunas empresas fijaban su precio antes de observar los choques agregados.
    \item Para estudiar inflación necesitamos describir cómo evoluciona el nivel general de precios cuando sólo una fracción de las empresas puede reajustarlo.
    \item La relación resultante conecta la inflación presente con:
    \begin{itemize}
      \item el nivel de actividad respecto de la producción con precios flexibles; y
      \item la inflación futura anticipada.
    \end{itemize}
    \item Esta relación es la \cidehighlight{curva de Phillips Nuevo-Keynesiana}.
  \end{itemize}
\end{frame}

\begin{frame}{Fijación de precios a la Calvo}
  \begin{itemize}
    \setlength{\itemsep}{0.58em}
    \item En cada periodo, cada empresa recibe aleatoriamente la oportunidad de reajustar su precio.
    \item La probabilidad de reajuste es independiente del tiempo transcurrido desde el último cambio.
    \item La empresa que reajusta elige un precio \textit{forward-looking}: pondera los costos marginales y la demanda esperados mientras su precio permanezca fijo.
    \item Las empresas que no reajustan satisfacen la demanda al precio previamente elegido.
  \end{itemize}

  \begin{block}{Implicación}
    Un mayor nivel de actividad eleva los costos marginales y el precio óptimo; una mayor inflación futura esperada induce a reajustar más el precio hoy.
  \end{block}
\end{frame}

\begin{frame}{Curva de Phillips Nuevo-Keynesiana}
  Definamos la brecha de producto con la convención usual:
  \[
    x_t \equiv Y_t-Y_t^{m},
  \]
  donde \(Y_t^{m}\) es la producción que prevalecería con precios flexibles. Entonces,
  \[
    \boxed{\pi_t=\kappa x_t+\beta E_t\pi_{t+1}},
    \qquad \kappa>0,\quad 0<\beta<1.
  \]

  \begin{itemize}
    \setlength{\itemsep}{0.45em}
    \item Si \(Y_t>Y_t^m\), los costos marginales y la inflación aumentan.
    \item Si \(Y_t<Y_t^m\), la brecha es negativa y la inflación disminuye.
    \item Una mayor inflación futura esperada eleva la inflación presente.
  \end{itemize}

  \begin{alertblock}{Convención de Williamson}
    El libro denomina brecha recesiva a \(Y_t^m-Y_t=-x_t\). Ambas convenciones describen el mismo mecanismo.
  \end{alertblock}
\end{frame}

\begin{frame}{Relación de Fisher y tasa natural}
  \begin{itemize}
    \setlength{\itemsep}{0.62em}
    \item El banco central controla la tasa de interés nominal \(R_t\), pero las decisiones intertemporales dependen de la tasa real:
    \[
      \boxed{r_t=R_t-E_t\pi_{t+1}}.
    \]
    \item La \cidehighlight{tasa real natural} \(r_t^n\) es la tasa compatible con \(Y_t=Y_t^m\), es decir, con una brecha de producto igual a cero.
    \item Si \(r_t>r_t^n\), el gasto presente es insuficiente y \(x_t<0\).
    \item Si \(r_t<r_t^n\), el gasto presente es elevado y \(x_t>0\).
  \end{itemize}
\end{frame}

\section{Modelo Nuevo-Keynesiano básico}

\begin{frame}{El modelo Nuevo-Keynesiano básico}
  \centering
  \figura[6.25cm]{fig_15_1_modelo_nk.png}
  \fuentefigura{Williamson (2018), figura 15.1.}
\end{frame}

\begin{frame}{Determinación del equilibrio}
  \begin{enumerate}
    \setlength{\itemsep}{0.72em}
    \item El banco central fija la tasa nominal \(R_t\).
    \item Dada la inflación futura esperada, se determina la tasa real:
    \[
      r_t=R_t-E_t\pi_{t+1}.
    \]
    \item La tasa real determina el gasto deseado y el producto \(Y_t=Y^d(r_t)\).
    \item Dado \(Y_t\), la curva de Phillips determina la inflación presente:
    \[
      \pi_t=\kappa(Y_t-Y_t^m)+\beta E_t\pi_{t+1}.
    \]
  \end{enumerate}

  \begin{block}{Modelo básico}
    En esta primera versión, \(E_t\pi_{t+1}\) se trata como una variable exógena.
  \end{block}
\end{frame}

\begin{frame}{Mandato dual del banco central}
  \begin{itemize}
    \setlength{\itemsep}{0.55em}
    \item El banco central desea estabilizar la inflación alrededor de su meta \(\pi^*\).
    \item También desea cerrar la brecha de producto: \(Y_t=Y_t^m\).
    \item Una representación reducida de sus pérdidas es
    \[
      \mathcal{L}_t=(\pi_t-\pi^*)^2+\lambda x_t^2,
      \qquad \lambda>0.
    \]
    \item Si \(r_t=r_t^n\) y la inflación esperada es consistente con la meta, ambos objetivos pueden alcanzarse simultáneamente.
    \item Ante ciertos choques, el banco central enfrenta una disyuntiva entre inflación y actividad.
  \end{itemize}
\end{frame}

\begin{frame}{Caída de la tasa real natural}
  \centering
  \figura[6.25cm]{fig_15_2_tasa_natural.png}
  \fuentefigura{Williamson (2018), figura 15.2.}
\end{frame}

\begin{frame}{El dilema de política monetaria}
  \begin{itemize}
    \setlength{\itemsep}{0.62em}
    \item Una caída de \(r^n\) deja a la tasa real vigente por encima de su nivel natural.
    \item Si el banco central no modifica \(R\), mantiene la inflación en la meta, pero permite una brecha recesiva.
    \item Si reduce \(R\) hasta igualar \(r=r^n\), cierra la brecha, pero la inflación supera la meta.
    \item Si valora ambos objetivos, elige una respuesta intermedia: tolera desviaciones moderadas de inflación y producto.
  \end{itemize}

  \begin{block}{Idea central}
    El choque desplaza el nivel de actividad compatible con precios flexibles; la rigidez nominal impide alcanzar simultáneamente ambos objetivos.
  \end{block}
\end{frame}

\begin{frame}{Aumento de la inflación futura esperada}
  \centering
  \figura[6.25cm]{fig_15_3_inflacion_esperada.png}
  \fuentefigura{Williamson (2018), figura 15.3.}
\end{frame}

\begin{frame}{Respuesta ante mayores expectativas de inflación}
  \begin{itemize}
    \setlength{\itemsep}{0.58em}
    \item El aumento de \(E_t\pi_{t+1}\) desplaza la curva de Phillips hacia arriba.
    \item Si \(R_t\) aumenta uno a uno con la inflación esperada, \(r_t\) no cambia: el producto permanece constante, pero la inflación aumenta.
    \item Para mantener \(\pi_t=\pi^*\), el banco central debe aumentar \(R_t\) más que proporcionalmente, elevando \(r_t\) y reduciendo el producto.
    \item Una respuesta intermedia distribuye el costo entre inflación por encima de la meta y una brecha recesiva.
  \end{itemize}

  \begin{alertblock}{Principio de Taylor}
    La lógica convencional prescribe aumentar la tasa nominal más que uno a uno ante un aumento persistente de la inflación.
  \end{alertblock}
\end{frame}

\section{Evidencia y límite inferior}

\begin{frame}{La curva de Phillips en los datos}
  \begin{itemize}
    \setlength{\itemsep}{0.58em}
    \item Phillips documentó originalmente una relación negativa entre inflación salarial y desempleo.
    \item En el modelo moderno, la inflación depende de la brecha de actividad y de las expectativas de inflación.
    \item Por ello, una correlación simple entre inflación y desempleo no identifica una curva estructural estable.
    \item Cambios en expectativas, choques de oferta y modificaciones de la regla monetaria desplazan la relación observada.
  \end{itemize}

  \begin{block}{Predicción condicional}
    Manteniendo constantes las expectativas y los choques de costos, una mayor actividad económica debería asociarse con mayor inflación.
  \end{block}
\end{frame}

\begin{frame}{Una curva de Phillips para Estados Unidos}
  \centering
  \figura[6.10cm]{fig_15_4_phillips_1949_2016.png}
  \fuentefigura{Williamson (2018), figura 15.4.}
\end{frame}

\begin{frame}{La relación después de la Gran Recesión}
  \centering
  \figura[6.10cm]{fig_15_5_phillips_2009_2016.png}
  \fuentefigura{Williamson (2018), figura 15.5.}
\end{frame}

\begin{frame}{¿Qué aprendemos de la evidencia?}
  \begin{itemize}
    \setlength{\itemsep}{0.64em}
    \item La relación incondicional entre inflación y holgura económica es débil e inestable.
    \item Una curva de Phillips estructural no implica que todos los movimientos observados deban ubicarse sobre una misma curva.
    \item La pendiente estimada puede cambiar con el anclaje de expectativas, la frecuencia de reajuste de precios y la composición de los choques.
    \item La evidencia exige cautela: la curva de Phillips es un mecanismo del modelo, no una regularidad mecánica e invariable.
  \end{itemize}
\end{frame}

\begin{frame}{¿Por qué puede ser baja la tasa natural?}
  \begin{itemize}
    \setlength{\itemsep}{0.58em}
    \item \textbf{Exceso mundial de ahorro:} una mayor oferta de fondos prestables reduce la tasa real de equilibrio.
    \item \textbf{Estancamiento secular:} la escasez de oportunidades rentables reduce la demanda de inversión.
    \item \textbf{Fricciones financieras:} una mayor demanda de activos seguros y restricciones de crédito deprimen la tasa segura.
    \item En los tres casos, \(r_t^n\) puede caer a niveles muy bajos o incluso negativos.
  \end{itemize}

  \begin{block}{Problema de política}
    Si la tasa nominal no puede reducirse lo suficiente, la tasa real permanece por encima de su nivel natural.
  \end{block}
\end{frame}

\begin{frame}{Límite inferior de la tasa nominal}
  Si la tasa nominal satisface \(R_t\geq 0\), entonces
  \[
    r_t=R_t-E_t\pi_{t+1}\geq -E_t\pi_{t+1}.
  \]
  Cuando
  \[
    \boxed{r_t^n<-E_t\pi_{t+1}},
  \]
  ni siquiera \(R_t=0\) permite alcanzar la tasa real natural.

  \begin{itemize}
    \setlength{\itemsep}{0.50em}
    \item La tasa real resulta demasiado alta.
    \item El producto queda por debajo de \(Y_t^m\).
    \item La inflación queda por debajo de su meta.
    \item La política monetaria convencional queda restringida.
  \end{itemize}
\end{frame}

\begin{frame}{Trampa de liquidez}
  \centering
  \figura[6.25cm]{fig_15_6_zlb.png}
  \fuentefigura{Williamson (2018), figura 15.6.}
\end{frame}

\begin{frame}{Orientación futura}
  La \cidehighlight{orientación futura} (\textit{forward guidance}) es un compromiso sobre la política monetaria futura.

  \[
    E_t\pi_{t+1}\uparrow
    \quad\Longrightarrow\quad
    r_t=R_t-E_t\pi_{t+1}\downarrow
    \quad\Longrightarrow\quad
    Y_t\uparrow,\ \pi_t\uparrow.
  \]

  \begin{itemize}
    \setlength{\itemsep}{0.52em}
    \item En el límite inferior, \(R_t\) no puede caer, pero las expectativas de inflación sí pueden reducir la tasa real.
    \item El banco central promete mantener una política expansiva en el futuro, incluso después de que la economía se recupere.
    \item La promesa desplaza la curva de Phillips y estimula el gasto presente.
  \end{itemize}
\end{frame}

\begin{frame}{Orientación futura en el límite inferior}
  \centering
  \figura[6.25cm]{fig_15_7_forward_guidance.png}
  \fuentefigura{Williamson (2018), figura 15.7.}
\end{frame}

\begin{frame}{Credibilidad e inconsistencia temporal}
  \begin{itemize}
    \setlength{\itemsep}{0.62em}
    \item La orientación futura funciona sólo si modifica las expectativas privadas.
    \item Cuando llegue el futuro, el banco central tendrá incentivos para incumplir la promesa y evitar inflación excesiva.
    \item Si el público anticipa ese incumplimiento, la promesa no será creíble y no reducirá la tasa real presente.
    \item La comunicación debe ser clara y estar respaldada por mecanismos de compromiso.
  \end{itemize}

  \begin{alertblock}{Tensión fundamental}
    La política que conviene prometer hoy puede no ser la que convenga ejecutar mañana.
  \end{alertblock}
\end{frame}

\section{Expectativas racionales y neo-fisherismo}

\begin{frame}{Modelo NK con expectativas racionales}
  Supongamos ausencia de choques y parámetros constantes. La demanda intertemporal satisface
  \[
    Y_t-Y_{t+1}
    =-\frac{1}{d}\left(R_t-\pi_{t+1}-r^*\right),
    \qquad d>0.
  \]
  Williamson simplifica la curva de Phillips a
  \[
    \pi_t=a(Y_t-Y^m),\qquad a>0,
  \]
  y, en el periodo siguiente,
  \[
    \pi_{t+1}=a(Y_{t+1}-Y^m).
  \]

  \begin{block}{Expectativas racionales}
    Sin choques, los valores anticipados coinciden con los realizados: los agentes conocen el modelo y no cometen errores sistemáticos.
  \end{block}
\end{frame}

\begin{frame}{Ley de movimiento de la inflación}
  Al eliminar \(Y_t\) y \(Y_{t+1}\), obtenemos
  \[
    \boxed{
      \pi_{t+1}
      =
      \frac{a(R_t-r^*)}{a+d}
      +
      \frac{d}{a+d}\pi_t
    }.
  \]

  \begin{itemize}
    \setlength{\itemsep}{0.58em}
    \item La inflación presente y la tasa nominal determinan la inflación futura.
    \item Como \(0<d/(a+d)<1\), para una tasa nominal constante la inflación converge a un estado estacionario.
    \item La tasa nominal fija el estado estacionario, pero no determina por sí sola la inflación inicial.
  \end{itemize}

  \begin{alertblock}{Indeterminación}
    Existen múltiples trayectorias de inflación compatibles con una tasa nominal constante.
  \end{alertblock}
\end{frame}

\begin{frame}{Tasa nominal constante}
  \centering
  \figura[5.95cm]{fig_15_8_nkre_constante.png}
  \fuentefigura{Williamson (2018), figura 15.8.}
\end{frame}

\begin{frame}{Relación de Fisher de largo plazo}
  En estado estacionario, \(\pi_{t+1}=\pi_t=\pi\). Por lo tanto,
  \[
    \boxed{\pi=R-r^*}.
  \]

  \begin{itemize}
    \setlength{\itemsep}{0.64em}
    \item Si \(r^*\) permanece constante, un aumento permanente de \(R\) eleva la inflación de largo plazo uno a uno.
    \item Éste es el \cidehighlight{efecto Fisher de largo plazo}.
    \item En el modelo, la causalidad va de la tasa nominal fijada por el banco central hacia la inflación de estado estacionario.
    \item El resultado requiere interpretar el aumento de \(R\) como un cambio permanente y creíble del régimen monetario.
  \end{itemize}
\end{frame}

\begin{frame}{Aumento permanente de la tasa nominal}
  \centering
  \figura[5.85cm]{fig_15_9_alza_tasa_nkre.png}
  \fuentefigura{Williamson (2018), figura 15.9.}
\end{frame}

\begin{frame}{Dinámica de la inflación}
  \centering
  \figura[5.85cm]{fig_15_10_dinamica_inflacion.png}
  \fuentefigura{Williamson (2018), figura 15.10.}
\end{frame}

\begin{frame}{La intuición neo-fisheriana}
  \begin{itemize}
    \setlength{\itemsep}{0.58em}
    \item Una tasa nominal permanentemente mayor debe ser compatible, en el largo plazo, con una inflación mayor si \(r^*\) no cambia.
    \item Con expectativas racionales, los agentes anticipan la nueva trayectoria de inflación y ajustan sus decisiones intertemporales desde hoy.
    \item En el modelo NKRE de Williamson, la inflación no cae inicialmente: converge hacia el nuevo estado estacionario más alto.
    \item Esto contrasta con el experimento convencional de un aumento transitorio e inesperado de la tasa, que puede ser contractivo y desinflacionario.
  \end{itemize}

  \begin{alertblock}{Advertencia}
    El resultado neo-fisheriano depende del carácter permanente, creíble y fiscalmente consistente del régimen; no es una correlación mecánica entre tasas e inflación.
  \end{alertblock}
\end{frame}

\section{Reglas de Taylor}

\begin{frame}{Regla de Taylor simplificada}
  Williamson considera
  \[
    \boxed{
      R_t=\max\left\{0,\,
      h\pi_t+(1-h)\pi^*+r^*
      \right\}.
    }
  \]

  \begin{itemize}
    \setlength{\itemsep}{0.56em}
    \item \(\pi^*\) es la meta de inflación.
    \item El operador \(\max\) incorpora el límite inferior \(R_t\geq 0\).
    \item Si \(h>1\), la regla satisface el \cidehighlight{principio de Taylor}: la tasa nominal responde más que uno a uno a la inflación.
    \item La intención es elevar la tasa real cuando la inflación aumenta.
  \end{itemize}
\end{frame}

\begin{frame}{Dos estados estacionarios}
  \centering
  \figura[5.95cm]{fig_15_11_regla_taylor.png}
  \fuentefigura{Williamson (2018), figura 15.11.}
\end{frame}

\begin{frame}{Los peligros de la regla de Taylor}
  \centering
  \figura[5.95cm]{fig_15_12_dinamica_taylor.png}
  \fuentefigura{Williamson (2018), figura 15.12.}
\end{frame}

\begin{frame}{¿Qué puede salir mal?}
  \begin{itemize}
    \setlength{\itemsep}{0.62em}
    \item La regla tiene un estado estacionario deseado: \(\pi=\pi^*\) y \(R=r^*+\pi^*\).
    \item También puede tener un estado estacionario no deseado: \(R=0\) y \(\pi=-r^*\).
    \item Bajo el principio de Taylor, el estado deseado puede ser inestable y múltiples trayectorias convergen al equilibrio de baja inflación.
    \item El banco central observa inflación insuficiente, desea reducir \(R\), pero ya se encuentra en el límite inferior.
  \end{itemize}

  \begin{alertblock}{Peligro}
    Una regla diseñada para estabilizar la inflación puede producir indeterminación y una trampa persistente de baja inflación.
  \end{alertblock}
\end{frame}

\begin{frame}{Una regla monetaria neo-fisheriana}
  Williamson propone una regla por tramos. Si \(\pi_t<\pi^*\),
  \[
    R_t=r^*+\frac{a+d}{a}\pi^*-\frac{d}{a}\pi_t.
  \]
  Si \(\pi_t\geq\pi^*\),
  \[
    R_t=r^*-\frac{d}{a}\pi^*
        +\frac{a+d}{a}E_t\pi_{t+1}.
  \]

  \begin{itemize}
    \setlength{\itemsep}{0.48em}
    \item Con inflación baja, la regla prescribe una tasa nominal más alta.
    \item Con presiones inflacionarias futuras, la tasa responde más que proporcionalmente.
    \item En el modelo, \(\pi_{t+1}=\pi^*\) después de cualquier inflación inicial inferior a la meta.
  \end{itemize}
\end{frame}

\begin{frame}{Inflación bajo la regla neo-fisheriana}
  \centering
  \figura[5.75cm]{fig_15_13_regla_neofisheriana.png}
  \fuentefigura{Williamson (2018), figura 15.13.}
\end{frame}

\begin{frame}{Dos visiones de la política monetaria}
  \small
  \renewcommand{\arraystretch}{1.30}
  \begin{tabularx}{\textwidth}{lYY}
    \toprule
      & \textbf{Visión convencional} & \textbf{Visión neo-fisheriana} \\
    \midrule
    Experimento
      & Cambio transitorio de la tasa nominal.
      & Cambio permanente y creíble del régimen. \\
    Mecanismo inicial
      & Mayor tasa real reduce gasto y actividad.
      & Expectativas ajustan la trayectoria de inflación. \\
    Inflación
      & Una mayor tasa nominal reduce la inflación en el corto plazo.
      & Una mayor tasa nominal eleva la inflación en el largo plazo y, en el NKRE, también su trayectoria. \\
    Riesgo
      & Recesión por endurecimiento excesivo.
      & Indeterminación, falta de credibilidad o inconsistencia fiscal. \\
    \bottomrule
  \end{tabularx}
\end{frame}

\begin{frame}{Conclusiones}
  \begin{itemize}
    \setlength{\itemsep}{0.58em}
    \item La curva de Phillips conecta inflación, brecha de producto y expectativas, pero su relación empírica incondicional es débil.
    \item Una tasa natural baja puede llevar a la economía al límite inferior y volver insuficiente la política convencional.
    \item La orientación futura intenta reducir la tasa real mediante expectativas, pero exige credibilidad y compromiso.
    \item Con expectativas racionales, el modelo NK genera la relación de Fisher de largo plazo y resultados neo-fisherianos.
    \item Las reglas de Taylor pueden admitir múltiples equilibrios; el diseño de política debe considerar expectativas, permanencia y consistencia del régimen.
  \end{itemize}
\end{frame}

\begin{frame}[plain]
  \centering
  \vfill
  {\LARGE\color{CIDEBlue}\bfseries Inflación, expectativas y política monetaria}
  \vspace{0.55cm}

  {\large La dinámica depende tanto de la acción presente del banco central\\
  como de las creencias sobre su conducta futura.}
  \vfill
\end{frame}

\end{document}

